% Introduction.  Opening follows the CH2/CH7 grammar: question, setting,
% object, results, organisation.  Three candidate advantages kept as content.

\section{Introduction}
\label{path:sec:introduction}

Under linear clearance the schedule by which an infected cell releases its
load is invisible to the basic reproduction number: replacing constant
budding by a single burst leaves $R_0$ exactly unchanged, a fact proved in
\ChPop{} (\cref{p:thm:R0inv}). The schedule is not invisible to
establishment, but that comparison is again a comparison of branching
structures at fixed mean. This chapter asks whether a depletable defence
changes either conclusion --- whether bursting can be an advantage, or a
disadvantage, once the host's capacity to kill is finite.

Two sites are treated. Inside the host cell, some leucocytes carry a finite
store of oxidative granules, spent as they act. Outside it, killer cells in
the extracellular matrix remove free agents on contact and are themselves
used up. In both cases the per-agent hazard depends on how many agents
arrive at once. Constant budding, as used in the basic model of viral
replication \cite{anderson1988role,perelson1996hiv,nowak1996population}, is
an approximation that ignores that dependence; single-cell observations of
HIV release are already intermittent rather than constant
\cite{ribeiro2010estimation,capistran2018extracellular,ciupe2017host}.
Eclipse phases and load-dependent production
\cite{nelson2004age,kitagawa2018pde,rong2013analysis} address a different
approximation, and are not pursued here.

Comparing bursting with budding is not new
\cite{pearson2011stochastic,komarova2007viral,bird2015escape,bird2015nonlytic,garoff1998virus}.
What is added is a quantification of \emph{flooding} --- that a burst-style
release is more likely to exhaust the local killing capacity. The word is
used below in two senses that must be kept apart --- a comparison of two
offspring laws at one delivery, and a comparison of the same yield split
across more deliveries --- and the two do not always agree.

\subsection{Three candidate advantages}
\label{path:sec:advantages}

Is there a fundamental advantage to bursting dynamics in pathogenesis? For
\ft{} and \yp{}, host cell lysis may be no more than the simplest available
means of release: bacteria are large in comparison with the protein complexes
that constitute a cell membrane, so it would be difficult to employ a
strategy in which some but not all of the pathogen were released. Viruses are
different. They are small enough to allow partial release of an
intracellular load, and for viruses such as HIV there is an explicit
advantage to not killing the cell, since its machinery continues to produce
virions. That an all-at-once strategy nonetheless persists is what needs
explaining, and the explanations on offer fall into three categories.

\begin{description}[leftmargin=1.4em,style=unboxed,itemsep=0.2em]
\item[Subductive proliferation.] Replication rate declines over time within
a host cell, perhaps following recognition, or as the resources available
there are used up.
\item[Limited killing.] The means of removing agents are finite on the
relevant timescales, so the effective death rate of invading pathogen falls
as the incident cohort is killed. This may happen inside the cell, where
some host cells carry oxidative granules of finite number, or outside it.
\item[Detection avoidance.] Continual budding may be detected more readily,
by immune cells that induce apoptosis, than release that is sporadic and
clustered, as HIV's is observed to be.
\end{description}

All three amount to catching some immune mechanism off guard. This chapter
pursues the second, because it is the one that can be settled by
calculation; the other two are named here and left. Subductive proliferation
would need no new machinery, its dynamics being the standard
time-inhomogeneous theory already available in \cref{m:sec:inhomogeneous},
but that theory does not by itself distinguish bursting from budding.
Detection avoidance would need a model of detection, which this chapter does
not have. \Cref{path:sec:limitations} returns to both.

Three results carry what follows. If a fixed count $\vartheta$ of released
agents is removed on each delivery, the expected surviving release of one
cell is $\mathcal{R}(\vartheta)=V_\infty a^{-\vartheta}$, where $V_\infty$
is the expected lifetime yield of an infected cell and $a$ is the larger
root of the characteristic quadratic taken from \ChCore{}. Every unit of
removal costs the burst one factor of $a$. In the replicator--suppressor
model of \cref{path:sec:intracellular}, survival is certain once the
incident cohort exceeds the granule store; on the matched diagonal the
infection fails with a probability that decays only polynomially in cohort
size, at exponent $\lambda/\varsigma$. And the flooding criterion of
\ChPop{} reverses under a removal threshold: bursting wins the matched-mean
comparison exactly when $L<1$, the complement of \cref{p:thm:flood}.

The derivation proceeds as follows.
\begin{itemize}[leftmargin=1.6em,itemsep=0.25em,topsep=0.35em]
\item A chained birth--death--catastrophe process with complete transfer
cannot prefer one release schedule, so a nonlinearity is forced into the
model.
\item The intracellular chain is solved by first-step analysis; the
certainty threshold and the diagonal law are closed.
\item The extracellular toll is reduced to an integer budget, and
memorylessness of the burst law gives $\mathcal{R}(\vartheta)$.
\item The same functional is compared across bursting and budding, and the
criterion reverses.
\end{itemize}

The intracellular process, its roots $a$ and $b$, its lifetime yield
$V_\infty$ and its burst-size law are taken from \ChCore{}
(\cref{bdc:lem:roots,bdc:eq:Vinf,dist:eq:burstlaw}); the single-cell
comparison of the two release strategies is established in \ChDist{}
(\cref{dist:sec:budding}); and the population comparison, the flooding
theorem and the catalogue of intermediate release mechanisms are proved in
\ChPop{} (\cref{p:thm:flood,p:sec:reach}). New here are the two depletion
models, the closed form $\mathcal{R}(\vartheta)=V_\infty a^{-\vartheta}$,
and the reversal. \Cref{path:sec:schedule} forces the nonlinearity;
\cref{path:sec:intracellular,path:sec:removal} supply it, one site each;
\cref{path:sec:selection} places the result against the flooding theorem and
closes on \yp{}. \Cref{path:app:verification} records the numerical checks.

\begin{table}[htbp]
\centering
\small
\begin{tabular}{@{}ll@{}}
\toprule
Symbol & Meaning \\
\midrule
\multicolumn{2}{@{}l}{\itshape Inherited, and used here without redefinition} \\
$\lambda,\mu,\delta$ & birth, death and catastrophe rates of one replicator \\
$a,b$ & roots of the characteristic quadratic, $b<1<a$ \\
$L=a(1-b)$ & the flooding parameter of \cref{p:thm:flood} \\
$W_t$ & agents released by one infected cell up to time $t$ \\
$V_\infty$ & expected lifetime release of one infected cell, $\lim_t\ex{W_t}$ \\
$r$ & burst size: the load released at rupture \\
$q$ & probability that a released agent infects a new cell \\
\addlinespace[2pt]
\multicolumn{2}{@{}l}{\itshape Particular to this chapter} \\
$\vartheta$ & fixed removal budget, an integer count of killed agents \\
$\vartheta^{\ast}$ & the budget at which $\mathcal{R}=1$, $\log V_\infty/\log a$ \\
$\mathcal{R}(\vartheta)$ & expected surviving release per cell under budget $\vartheta$ \\
$X_t,\;Q_t$ & intracellular replicators, and remaining granules \\
$\varsigma$ & suppression rate per replicator--granule pair \\
$\pi_{i,j}$ & probability of extinction from $i$ replicators, $j$ granules \\
$N$ & granule store of a fresh host cell \\
\bottomrule
\end{tabular}
\caption{Notation. The upper block is taken from \ChCore{} and \ChPop{}; the
lower block carries meanings specific to this chapter, and those letters
were chosen to avoid collisions with the ones those meanings would otherwise
have taken.}
\label{path:tab:notation}
\end{table}

\FloatBarrier
